% Sec. XI Conclusion (about 0.2 page): what it gives, what it showed, what is next.
\section{Conclusion}
\label{sec:conclusion}
SCooP-iT captures a real heterogeneous indoor team, sensing and
communication alike, on one clock, with measured offsets below
\SI{1}{\milli\second}, and in one metric frame, placed to
\SI{10}{\milli\metre} against held-out boards without motion capture or
GNSS. The capture serves three uses, each checked against reality: the
recordings as a test bench into which any collaborative method is
replayed against the link the agents had \todo{validated to within one
receiver frame}; SCooP-iT Gym, a closed-loop environment whose world,
sensors, link, and conditions are built from the recordings; and the
testbed itself, on which a policy trained in the Gym is deployed and
its outcome compared with the Gym's prediction. The real world showed
what simulation hides: under contention the shared network moves raw
clouds past one LiDAR frame while its own latency looks healthy, which
no synthetic channel of prior benchmarks reproduces. Next steps are more
agents and sites, other radios, multi-antenna CSI, a wired CSI host, and
Gym tasks in which agents act on the real testbed in closed loop. The
sequences recorded on SCooP-iT, their labels, and their benchmark are
released as the SCooP dataset \cite{scoop_dataset}.
